\documentclass[12pt,letterpaper]{article}
\usepackage{graphicx,textcomp}
\usepackage{natbib}
\usepackage{setspace}
\usepackage{fullpage}
\usepackage{color}
\usepackage[reqno]{amsmath}
\usepackage{amsthm}
\usepackage{fancyvrb}
\usepackage{amssymb,enumerate}
\usepackage[all]{xy}
\usepackage{endnotes}
\usepackage{lscape}
\newtheorem{com}{Comment}
\usepackage{float}
\usepackage{hyperref}
\newtheorem{lem} {Lemma}
\newtheorem{prop}{Proposition}
\newtheorem{thm}{Theorem}
\newtheorem{defn}{Definition}
\newtheorem{cor}{Corollary}
\newtheorem{obs}{Observation}
\usepackage[compact]{titlesec}
\usepackage{dcolumn}
\usepackage{tikz}
\usetikzlibrary{arrows}
\usepackage{multirow}
\usepackage{xcolor}
\newcolumntype{.}{D{.}{.}{-1}}
\newcolumntype{d}[1]{D{.}{.}{#1}}
\definecolor{light-gray}{gray}{0.65}
\usepackage{url}
\usepackage{listings}
\usepackage{color}

\definecolor{codegreen}{rgb}{0,0.6,0}
\definecolor{codegray}{rgb}{0.5,0.5,0.5}
\definecolor{codepurple}{rgb}{0.58,0,0.82}
\definecolor{backcolour}{rgb}{0.95,0.95,0.92}

\lstdefinestyle{mystyle}{
	backgroundcolor=\color{backcolour},
	commentstyle=\color{codegreen},
	keywordstyle=\color{magenta},
	numberstyle=\tiny\color{codegray},
	stringstyle=\color{codepurple},
	basicstyle=\footnotesize,
	breakatwhitespace=false,
	breaklines=true,
	captionpos=b,
	keepspaces=true,
	numbers=left,
	numbersep=5pt,
	showspaces=false,
	showstringspaces=false,
	showtabs=false,
	tabsize=2
}
\lstset{style=mystyle}
\newcommand{\Sref}[1]{Section~\ref{#1}}
\newtheorem{hyp}{Hypothesis}

\title{Problem Set 4}
\date{Due: April 10, 2026}
\author{Applied Stats II}


\begin{document}
	\maketitle
	\section*{Instructions}
	\begin{itemize}
	\item Please show your work! You may lose points by simply writing in the answer. If the problem requires you to execute commands in \texttt{R}, please include the code you used to get your answers. Please also include the \texttt{.R} file that contains your code. If you are not sure if work needs to be shown for a particular problem, please ask.
	\item Your homework should be submitted electronically on GitHub in \texttt{.pdf} form.
	\item This problem set is due before 23:59 on Friday April 10, 2026. No late assignments will be accepted.

	\end{itemize}

	\vspace{.25cm}
\section*{Question 1}
\vspace{.25cm}
\noindent We're interested in modeling the historical causes of child mortality. We have data from 26855 children born in Skellefteå, Sweden from 1850 to 1884. Using the ``child'' dataset in the \texttt{eha} library, fit a Cox Proportional Hazard model using mother's age and infant's gender as covariates. Present and interpret the output.

\vspace{.5cm}

\lstinputlisting[language=R, firstline=42, lastline=52]{PS04_SN.R}
\vspace{.5cm}


% Table created by stargazer v.5.2.3 by Marek Hlavac, Social Policy Institute. E-mail: marek.hlavac at gmail.com
% Date and time: Sat, Apr 04, 2026 - 17:26:18
\begin{table}[!htbp] \centering 
  \caption{Cox Proportional Hazard Model: Child Mortality} 
  \label{} 
\begin{tabular}{@{\extracolsep{5pt}}lc} 
\\[-1.8ex]\hline 
\hline \\[-1.8ex] 
 & \multicolumn{1}{c}{\textit{Dependent variable:}} \\ 
\cline{2-2} 
\\[-1.8ex] & enter \\ 
\hline \\[-1.8ex] 
 m.age & 0.008$^{***}$ \\ 
  & (0.002) \\ 
  & \\ 
 sexfemale & $-$0.082$^{***}$ \\ 
  & (0.027) \\ 
  & \\ 
\hline \\[-1.8ex] 
Observations & 26,574 \\ 
R$^{2}$ & 0.001 \\ 
Max. Possible R$^{2}$ & 0.986 \\ 
Log Likelihood & $-$56,503.480 \\ 
Wald Test & 22.520$^{***}$ (df = 2) \\ 
LR Test & 22.518$^{***}$ (df = 2) \\ 
Score (Logrank) Test & 22.530$^{***}$ (df = 2) \\ 
\hline 
\hline \\[-1.8ex] 
\textit{Note:}  & \multicolumn{1}{r}{$^{*}$p$<$0.1; $^{**}$p$<$0.05; $^{***}$p$<$0.01} \\ 
\end{tabular} 
\end{table}  


\subsection*{Interpretation}

The Cox Proportional Hazard model estimates the effect of mother's age (\texttt{m.age}) and infant's gender (\texttt{sex}) on the hazard of child mortality, using 26,574 observations with 5,616 death events. The model is jointly significant (likelihood ratio test: $\chi^2 = 22.52$, $df = 2$, $p < 0.001$).

\textbf{Mother's age (\texttt{m.age}):} The coefficient is 0.0076 ($p < 0.001$), with a hazard ratio of 1.008 (95\% CI: [1.003, 1.012]). For each additional year of the mother's age, the hazard of child mortality increases by approximately 0.8\%. While statistically significant, the effect is substantively small.

\textbf{Infant's sex (\texttt{sexfemale}):} The coefficient is $-0.082$ ($p = 0.002$), with a hazard ratio of 0.921 (95\% CI: [0.874, 0.971]). Female infants have approximately 7.9\% lower hazard of mortality compared to male infants, consistent with the well-documented biological survival advantage for female infants.

The Cox model makes no assumption about the baseline hazard function but assumes that the covariate effects are constant over time. The concordance of 0.519 shows the model is barely better than chance at
distinguishing higher-risk from lower-risk children, suggesting that while mother's age and infant's sex are statistically significant predictors, they explain only a small portion of the variation in child mortality.

\newpage
\section*{Question 2}
\vspace{.25cm}
\noindent We want to estimate how the amount of damage caused by a disaster relates to (1) whether disaster relief is provided and (2) the amount of relief that is provided conditional on a donation occurring when a disaster hits a given country. Estimate a Heckman selection model using the \texttt{disasters\_response} data in which \texttt{binContribution} is the outcome for the selection equation, and \texttt{originalContributionMillionUSDLogged} is the outcome for the second equation. Use \texttt{occurrences + deathsEM + normalizedDamageEMLogged} as your input variables for both equations. Present and interpret the output.

\vspace{.5cm}

\lstinputlisting[language=R, firstline=63, lastline=83]{PS04_SN.R}
\vspace{.5cm}


% Table created by stargazer v.5.2.3 by Marek Hlavac, Social Policy Institute. E-mail: marek.hlavac at gmail.com
% Date and time: Sat, Apr 04, 2026 - 17:26:20
\begin{table}[!htbp] \centering 
  \caption{Heckman Selection Equation (Probit)} 
  \label{} 
\begin{tabular}{@{\extracolsep{5pt}} ccccc} 
\\[-1.8ex]\hline 
\hline \\[-1.8ex] 
 & Estimate & Std. Error & t value & Pr(\textgreater \textbar t\textbar ) \\ 
\hline \\[-1.8ex] 
(Intercept) & $$-$1.302$ & $0.018$ & $$-$72.276$ & $0$ \\ 
occurrences & $0.019$ & $0.002$ & $11.573$ & $0$ \\ 
deathsEM & $0.012$ & $0.001$ & $15.981$ & $0$ \\ 
normalizedDamageEMLogged & $0.021$ & $0.001$ & $23.360$ & $0$ \\ 
\hline \\[-1.8ex] 
\end{tabular} 
\end{table}  

\vspace{.5cm}

% Table created by stargazer v.5.2.3 by Marek Hlavac, Social Policy Institute. E-mail: marek.hlavac at gmail.com
% Date and time: Sat, Apr 04, 2026 - 17:26:20
\begin{table}[!htbp] \centering 
  \caption{Heckman Outcome Equation (OLS with correction)} 
  \label{} 
\begin{tabular}{@{\extracolsep{5pt}} ccccc} 
\\[-1.8ex]\hline 
\hline \\[-1.8ex] 
 & Estimate & Std. Error & t value & Pr(\textgreater \textbar t\textbar ) \\ 
\hline \\[-1.8ex] 
(Intercept) & $9.766$ & $6.044$ & $1.616$ & $0.106$ \\ 
occurrences & $$-$0.086$ & $0.055$ & $$-$1.572$ & $0.116$ \\ 
deathsEM & $$-$0.033$ & $0.026$ & $$-$1.284$ & $0.199$ \\ 
normalizedDamageEMLogged & $$-$0.114$ & $0.062$ & $$-$1.833$ & $0.067$ \\ 
\hline \\[-1.8ex] 
\end{tabular} 
\end{table}  

\vspace{.5cm}

% Table created by stargazer v.5.2.3 by Marek Hlavac, Social Policy Institute. E-mail: marek.hlavac at gmail.com
% Date and time: Sat, Apr 04, 2026 - 17:26:20
\begin{table}[!htbp] \centering 
  \caption{Heckman Model Diagnostics} 
  \label{} 
\begin{tabular}{@{\extracolsep{5pt}} ccccc} 
\\[-1.8ex]\hline 
\hline \\[-1.8ex] 
 & Estimate & Std. Error & t value & Pr(\textgreater \textbar t\textbar ) \\ 
\hline \\[-1.8ex] 
invMillsRatio & $$-$6.313$ & $3.420$ & $$-$1.846$ & $0.065$ \\ 
sigma & $5.992$ & $$ & $$ & $$ \\ 
rho & $$-$1.054$ & $$ & $$ & $$ \\ 
\hline \\[-1.8ex] 
\end{tabular} 
\end{table}  


\subsection*{Interpretation}

The Heckman selection model addresses sample selection bias: we only observe the amount of disaster relief (\texttt{originalContributionMillionUSDLogged}) when a contribution actually occurs (\texttt{binContribution} = 1). Of the 49,096 observations, 45,848 are censored (no contribution) and 3,248 are observed. Estimating the outcome equation on the selected sample alone would yield biased coefficients if the factors driving selection are correlated with the error in the contribution amount equation.

\textbf{Selection equation (Probit):} All three predictors are highly significant ($p < 0.001$). The number of disaster occurrences ($\hat{\beta} = 0.019$), deaths ($\hat{\beta} = 0.011$), and logged normalized damage ($\hat{\beta} = 0.021$) all positively predict the probability of receiving a contribution. More frequent, deadlier, and more damaging disasters are more likely to receive relief. The intercept is $-1.302$, indicating a low baseline probability of receiving a contribution.

\textbf{Outcome equation:} In contrast to the selection equation, none of the predictors are significant at the 5\% level in the outcome equation. The coefficients for occurrences ($\hat{\beta} = -0.086$, $p = 0.116$), deaths ($\hat{\beta} = -0.033$, $p = 0.199$), and logged normalized damage ($\hat{\beta} = -0.114$, $p = 0.067$) are all negative but not statistically significant. This suggests that while disaster severity drives \textit{whether} relief is given, it does not significantly predict \textit{how much} is given once a donation occurs. The adjusted $R^2$ is 0.053, indicating low explanatory power.

\textbf{Inverse Mills ratio ($\hat{\lambda}$):} The estimate is $-6.313$ ($p = 0.065$), marginally significant. The estimated $\hat{\rho} = -1.053$ (at the boundary, suggesting model instability) and $\hat{\sigma} = 5.992$. The marginal significance of $\hat{\lambda}$ provides weak evidence of selection bias.  This suggests unobservable factors that increase the likelihood of a contribution may be negatively associated with the contribution amount. Additionally, the boundary estimate of $\hat{\rho}$ suggests the two-step estimator may be performing poorly, and an alternative tool should be considered.


\vspace{1cm}
\section*{AI Usage Disclosure}
Claude Code (Anthropic, Claude Opus 4.6) was used to assist with R code development, \LaTeX{} formatting, and interpretation drafting for this problem set. Models used: Claude Opus 4.6 (1M context). All output was reviewed, verified, and edited by the author.

\end{document}
